\begin{lemma}[The scale-invariant Sobolev inequality at exponent two on a ball]
Let $n\ge2$ and suppose the Sobolev--Poincar\'e inequality on balls holds in
dimension $n$ (\noderef{CitedSobolevPoincareBall}).  Then there are
$\chi>1$ and $C>0$, depending only on $n$ and on the constant of that
inequality, such that for every $x_0\in\R^n$, every $\rho>0$ and every
$\varphi\in W^{1,2}(B_\rho(x_0))$ with weak gradient $g_\varphi$
(\noderef{MemW1pOn}) having compact support with
$\operatorname{supp}\varphi\subseteq B_\rho(x_0)$,
\[
\Bigl(\frac1{|B_\rho|}\int_{B_\rho}|\varphi|^{2\chi}dx\Bigr)^{\frac1{2\chi}}
\ \le\ C\,\rho\,
\Bigl(\frac1{|B_\rho|}\int_{B_\rho}|g_\varphi|^{2}dx\Bigr)^{\frac12},
\qquad B_\rho=B_\rho(x_0).
\]
One may take $\chi=\frac n{n-2}$ when $n\ge3$ and $\chi=2$ when $n=2$.

Both sides are averages, and the single factor $\rho$ makes the inequality
invariant under the dilations $x\mapsto x_0+\rho y$, which is the form the
Moser iteration over shrinking concentric balls consumes: no volume factors
have to be carried from one radius to the next.  The single parameter $\chi$ is
what lets one statement serve both $n\ge3$, where the Sobolev exponent is
$2n/(n-2)=2\chi$, and $n=2$, where there is no Sobolev exponent and any
$\chi>1$ will do.

This is not the cited result (C6).  \noderef{CitedSobolevPoincareBall} is the
$W^{1,1}\to L^{n/(n-1)}$ Poincar\'e--Sobolev inequality \emph{with the mean
subtracted}, for arbitrary $W^{1,1}$ functions on a ball; what is asserted here
is an $L^2\to L^{2\chi}$ inequality with \emph{no} mean, for functions with
compact support in the ball.  Since the Note for the formalization closes its
list of citable results with ``Everything else must be proved'', the passage
from the one to the other is proved below rather than assumed.
\end{lemma}
\begin{proof}
Throughout, $B_\sigma$ abbreviates $B_\sigma(x_0)$, and $\omega_n=|B_1(0)|$, so
that $|B_\sigma|=\omega_n\sigma^n$ by the Preamble item on the volume of a ball;
note $0<\omega_n<\infty$, so that $0<|B_\sigma|<\infty$ for every $\sigma>0$ and
the reciprocals $|B_\rho|^{-1}$ appearing in the statement are the honest ones.
Let $C_6=C_6(n)>0$ be the constant supplied by
\noderef{CitedSobolevPoincareBall}.  Set
\[
\chi=\frac n{n-2}\ \text{ and }\ \kappa=\frac{2(n-1)}{n-2}\quad(n\ge3),
\qquad
\chi=2\ \text{ and }\ \kappa=2\quad(n=2),
\]
and in both cases put $q=2\chi$, so $q=\frac{2n}{n-2}$ for $n\ge3$ and $q=4$
for $n=2$.  Then $\chi>1$: for $n\ge3$ because $n>n-2>0$, and for $n=2$
trivially.  Also $\kappa>1$ in both cases.

\emph{Step 1: the two exponent identities, for $n\ge3$.}  We record
\[
\kappa\cdot\frac n{n-1}=q
\qquad\text{and}\qquad
2(\kappa-1)=q
\qquad\text{and}\qquad
\frac{n-1}n-\frac12=\frac1q .
\tag{1}
\]
Indeed $\kappa\frac n{n-1}=\frac{2(n-1)}{n-2}\cdot\frac n{n-1}=\frac{2n}{n-2}=q$;
$2(\kappa-1)=2\bigl(\frac{2(n-1)}{n-2}-1\bigr)
=2\cdot\frac{2n-2-n+2}{n-2}=\frac{2n}{n-2}=q$; and
$\frac{n-1}n-\frac12=\frac{2(n-1)-n}{2n}=\frac{n-2}{2n}=\frac1q$.

\emph{Step 2: removing the mean, for a function vanishing on
$B_{2\rho}\setminus B_\rho$.}  Let $B'=B_{2\rho}$ and let $w\in L^1(B')$ with
$w=0$ almost everywhere on $B'\setminus B_\rho$; write
$w_{B'}=|B'|^{-1}\int_{B'}w$ and $s=\frac n{n-1}$.  We claim
\[
\Bigl(\int_{B'}|w|^{s}\Bigr)^{1/s}
\ \le\ c_n\Bigl(\int_{B'}|w-w_{B'}|^{s}\Bigr)^{1/s},
\qquad c_n=1+\bigl(1-2^{-n}\bigr)^{-1/s} .
\tag{2}
\]
By the Preamble item on the volume of a ball, $|B'|=2^n|B_\rho|$, so
$|B'\setminus B_\rho|=|B'|-|B_\rho|=(1-2^{-n})|B'|$.  On $B'\setminus B_\rho$ we
have $|w-w_{B'}|=|w_{B'}|$ almost everywhere, whence
\[
\int_{B'}|w-w_{B'}|^{s}
\ \ge\ \int_{B'\setminus B_\rho}|w_{B'}|^{s}
=(1-2^{-n})|B'|\,|w_{B'}|^{s},
\]
that is $\bigl(\int_{B'}|w-w_{B'}|^{s}\bigr)^{1/s}\ge
\bigl((1-2^{-n})|B'|\bigr)^{1/s}|w_{B'}|$.  On the other hand the $L^{s}$
triangle inequality (Minkowski, Preamble) applied to
$w=(w-w_{B'})+w_{B'}$ gives
\[
\Bigl(\int_{B'}|w|^{s}\Bigr)^{1/s}
\le\Bigl(\int_{B'}|w-w_{B'}|^{s}\Bigr)^{1/s}+|B'|^{1/s}\,|w_{B'}| ,
\]
the second term being the $L^{s}(B')$ norm of the constant $w_{B'}$.  Combining
the two displays and using
$|B'|^{1/s}\bigl((1-2^{-n})|B'|\bigr)^{-1/s}=(1-2^{-n})^{-1/s}$, which does not
depend on $\rho$, gives (2).  (Note $w_{B'}$ is finite because $w\in L^1(B')$.)

\emph{Step 3: the inequality for a smooth $\varphi$, case $n\ge3$.}  Let
$\varphi\in C^\infty(\R^n)$ have compact support with
$\operatorname{supp}\varphi\subseteq B_\rho$, and write $\nabla\varphi$ for its
classical gradient.  We prove
\[
\Bigl(\int_{B_\rho}|\varphi|^{q}\Bigr)^{1/q}
\ \le\ C_1\Bigl(\int_{B_\rho}|\nabla\varphi|^2\Bigr)^{1/2},
\qquad C_1=c_nC_6\kappa .
\tag{3}
\]
Put $B'=B_{2\rho}$, an open set of finite measure containing
$\operatorname{supp}\varphi$, and let $\Lambda=\sup_{\R^n}|\varphi|<\infty$
($\varphi$ is continuous with compact support).

\emph{(3a) $\varphi\in W^{1,1}(B')$ with weak gradient $\nabla\varphi$.}  This is
\noderef{SmoothCompactSupportMemW1p} at the exponent $r=1$: its hypotheses are
$n\ge1$, $1\le1$, openness of $B'$, smoothness of $\varphi$, compactness of its
support and $\operatorname{supp}\varphi\subseteq B'$, all of which hold.

\emph{(3b) $\psi=|\varphi|^{\kappa}$ lies in $W^{1,1}(B')$ with weak gradient
$\kappa|\varphi|^{\kappa-1}\operatorname{sgn}(\varphi)\nabla\varphi$.}  Let
$F(t)=|t|^{\kappa}$.  Since $\kappa>1$, $F$ is differentiable with
$F'(t)=\kappa|t|^{\kappa-1}\operatorname{sgn}(t)$, and $F'$ is continuous,
including at $t=0$ where $|F'(t)|\le\kappa|t|^{\kappa-1}\to0$; so $F$ is of
class $C^1$.  Its derivative is however unbounded on $\R$, which
\noderef{WeakGradientC1Comp} does not allow, so replace $F$ by the function
$\widetilde F$ equal to $F$ on $[-\Lambda,\Lambda]$ and extended affinely
outside, i.e.\ $\widetilde F(t)=F(\Lambda)+F'(\Lambda)(t-\Lambda)$ for
$t>\Lambda$ and $\widetilde F(t)=F(-\Lambda)+F'(-\Lambda)(t+\Lambda)$ for
$t<-\Lambda$.  Then $\widetilde F$ is $C^1$ on $\R$, the value and the
derivative matching at $\pm\Lambda$ by construction, and
$|\widetilde F'|\le L:=\kappa\Lambda^{\kappa-1}$ everywhere, because $|F'|$ is
nondecreasing in $|t|$.  Since $|\varphi|\le\Lambda$ we have
$\widetilde F\circ\varphi=F\circ\varphi=\psi$ and
$\widetilde F'\circ\varphi=F'\circ\varphi$.  Now
\noderef{WeakGradientC1Comp} applies with $q=1$, this $L$, the open set $B'$ of
finite measure, the pair $(\varphi,\nabla\varphi)$ of (3a) and the function
$\widetilde F$, and gives $\psi\in W^{1,1}(B')$ with weak gradient
$x\mapsto\widetilde F'(\varphi(x))\nabla\varphi(x)
=\kappa|\varphi(x)|^{\kappa-1}\operatorname{sgn}(\varphi(x))\nabla\varphi(x)$.
In particular, pointwise,
\[
|\nabla\psi|=\kappa\,|\varphi|^{\kappa-1}\,|\nabla\varphi| .
\tag{4}
\]

\emph{(3c) Applying (C6) on $B'$.}  By (3b) the pair $(\psi,\nabla\psi)$
satisfies the hypotheses of \noderef{CitedSobolevPoincareBall} on the ball
$B'=B_{2\rho}(x_0)$, whose radius $2\rho$ is positive; with $s=\frac n{n-1}$ its
conclusion reads
\[
\Bigl(\int_{B'}|\psi-\psi_{B'}|^{s}\Bigr)^{1/s}\le C_6\int_{B'}|\nabla\psi| .
\]
The function $\psi=|\varphi|^{\kappa}$ vanishes on $B'\setminus B_\rho$ and lies
in $L^1(B')$ (it is bounded by $\Lambda^{\kappa}$ and supported in a set of
finite measure), so Step 2 applies and gives
$\bigl(\int_{B'}\psi^{s}\bigr)^{1/s}\le c_nC_6\int_{B'}|\nabla\psi|$.
Inserting (4) and applying H\"older with the conjugate exponents $2$ and $2$ on
$B'$ (Preamble),
\[
\Bigl(\int_{B'}|\varphi|^{\kappa s}\Bigr)^{1/s}
\le c_nC_6\,\kappa\int_{B'}|\varphi|^{\kappa-1}|\nabla\varphi|
\le c_nC_6\,\kappa
\Bigl(\int_{B'}|\varphi|^{2(\kappa-1)}\Bigr)^{1/2}
\Bigl(\int_{B'}|\nabla\varphi|^{2}\Bigr)^{1/2}.
\]
By the first two identities of (1), $\kappa s=q$ and $2(\kappa-1)=q$, so with
$I=\int_{B'}|\varphi|^{q}$ this says
\[
I^{1/s}\ \le\ c_nC_6\kappa\,I^{1/2}\Bigl(\int_{B'}|\nabla\varphi|^2\Bigr)^{1/2}.
\]
Now $I<\infty$, because $|\varphi|^q\le\Lambda^q$ and $\varphi$ vanishes off a
set of finite measure.  If $I=0$ then $|\varphi|^q=0$ almost everywhere on
$B'$, so the left-hand side of (3) is $0$ and (3) holds, its right-hand side
being nonnegative.  If $I>0$, divide by $I^{1/2}$ and use the third identity of
(1), $\frac1s-\frac12=\frac1q$:
\[
I^{1/q}\le c_nC_6\kappa\Bigl(\int_{B'}|\nabla\varphi|^2\Bigr)^{1/2}.
\]
Finally $\varphi$ and $\nabla\varphi$ vanish off
$\operatorname{supp}\varphi\subseteq B_\rho$ --- for $\nabla\varphi$ because
$\varphi$ vanishes identically on the open complement of its support and a
function constant on an open set has vanishing derivative there --- so each
integral over $B'$ equals the corresponding integral over $B_\rho$.  This is
(3).

\emph{Step 4: the inequality for a smooth $\varphi$, case $n=2$.}  Here
$s=\frac n{n-1}=2$, $\kappa=2$ and $q=4$, and we prove
\[
\Bigl(\int_{B_\rho}|\varphi|^{4}\Bigr)^{1/4}
\ \le\ C_1\,\rho^{1/2}\Bigl(\int_{B_\rho}|\nabla\varphi|^2\Bigr)^{1/2},
\qquad C_1=\bigl(2c_2^2C_6^2\,\omega_2^{1/2}\bigr)^{1/2}.
\tag{5}
\]
Let $\varphi$, $B'$, $\Lambda$ be as in Step 3; (3a) and (3b) apply verbatim
with $\kappa=2$ (and there $F(t)=t^2$ is already $C^\infty$, so only the
truncation to bound $F'$ is needed), giving $\psi=\varphi^2\in W^{1,1}(B')$
with $|\nabla\psi|=2|\varphi||\nabla\varphi|$.

\emph{(4a) The Ladyzhenskaya step.}  Exactly as in (3c), with $s=2$:
\[
\Bigl(\int_{B'}|\varphi|^{4}\Bigr)^{1/2}
=\Bigl(\int_{B'}\psi^{2}\Bigr)^{1/2}
\le2c_2C_6\int_{B'}|\varphi||\nabla\varphi|
\le2c_2C_6\Bigl(\int_{B'}|\varphi|^{2}\Bigr)^{1/2}
\Bigl(\int_{B'}|\nabla\varphi|^{2}\Bigr)^{1/2},
\]
the last step H\"older with exponents $2$ and $2$ (Preamble).

\emph{(4b) A Poincar\'e step.}  Apply \noderef{CitedSobolevPoincareBall} on
$B'$ to the pair $(\varphi,\nabla\varphi)$ of (3a), again with $s=2$:
$\bigl(\int_{B'}|\varphi-\varphi_{B'}|^{2}\bigr)^{1/2}\le
C_6\int_{B'}|\nabla\varphi|$.  The function $\varphi$ vanishes on
$B'\setminus B_\rho$ and lies in $L^1(B')$, so Step 2 --- whose proof used only
that $w$ vanishes there, not that it is nonnegative --- gives
$\bigl(\int_{B'}|\varphi|^{2}\bigr)^{1/2}\le c_2C_6\int_{B'}|\nabla\varphi|$.
Since $\nabla\varphi$ vanishes off $B_\rho$, H\"older on the finite-measure set
$B_\rho$ (Preamble) bounds
$\int_{B'}|\nabla\varphi|=\int_{B_\rho}|\nabla\varphi|
\le|B_\rho|^{1/2}\bigl(\int_{B_\rho}|\nabla\varphi|^2\bigr)^{1/2}$, and
$|B_\rho|^{1/2}=(\omega_2\rho^2)^{1/2}=\omega_2^{1/2}\rho$ by the Preamble item
on the volume of a ball.  Hence
\[
\Bigl(\int_{B'}|\varphi|^{2}\Bigr)^{1/2}
\le c_2C_6\,\omega_2^{1/2}\,\rho\Bigl(\int_{B_\rho}|\nabla\varphi|^2\Bigr)^{1/2}.
\]

\emph{(4c) Combining.}  Substituting (4b) into (4a), and using that $\varphi$
and $\nabla\varphi$ vanish off $B_\rho$ so that the integrals over $B'$ are
integrals over $B_\rho$,
\[
\Bigl(\int_{B_\rho}|\varphi|^{4}\Bigr)^{1/2}
\le2c_2^2C_6^2\,\omega_2^{1/2}\,\rho\int_{B_\rho}|\nabla\varphi|^2 ,
\]
and taking square roots gives (5).

\emph{Step 5: from smooth $\varphi$ to the stated class.}  Let now
$(\varphi,g_\varphi)$ be as in the statement: $\varphi\in W^{1,2}(B_\rho)$ with
weak gradient $g_\varphi$ (\noderef{MemW1pOn}), compact support,
$\operatorname{supp}\varphi\subseteq B_\rho$.  Apply
\noderef{WeakGradientSmoothApprox} with this $n$, the exponent $q=2$, the open
set $B_\rho$ and the pair $(\varphi,g_\varphi)$; its hypotheses are $1\le n$,
$1\le2$, openness of $B_\rho$, the $W^{1,2}$ membership, compactness of the
support and $\operatorname{supp}\varphi\subseteq B_\rho$.  It produces a compact
$K\subseteq B_\rho$ with $\operatorname{supp}\varphi\subseteq K$ and functions
$\Phi_m\in C^\infty(\R^n)$, each with compact support contained in $K$, such
that
\[
\int_K|\Phi_m-\varphi|^2\to0
\qquad\text{and}\qquad
\int_K\bigl|D_i\Phi_m-(g_\varphi)_i\bigr|^2\to0
\quad(m\to\infty)
\tag{6}
\]
for every index $i$.

Each $\Phi_m$ is smooth with compact support contained in
$K\subseteq B_\rho$, so Step 3 (if $n\ge3$) resp.\ Step 4 (if $n=2$) applies to
it:
\[
\Bigl(\int_{B_\rho}|\Phi_m|^{q}\Bigr)^{1/q}
\le C_1\rho^{a}\Bigl(\int_{B_\rho}|\nabla\Phi_m|^2\Bigr)^{1/2},
\qquad a=0\ (n\ge3),\quad a=\tfrac12\ (n=2).
\tag{7}
\]
We pass to the limit in (7).

\emph{Right-hand side.}  All the integrands vanish off $K$, so the integrals
over $B_\rho$ are integrals over $K$.  By the reverse triangle inequality in
$\R^n$ and then the $L^2$ triangle inequality (Minkowski, Preamble),
\[
\Bigl|\Bigl(\int_K|\nabla\Phi_m|^2\Bigr)^{1/2}
-\Bigl(\int_K|g_\varphi|^2\Bigr)^{1/2}\Bigr|
\le\Bigl(\int_K|\nabla\Phi_m-g_\varphi|^2\Bigr)^{1/2}
\le\sum_{i=1}^n\Bigl(\int_K\bigl|D_i\Phi_m-(g_\varphi)_i\bigr|^2\Bigr)^{1/2},
\]
which tends to $0$ by (6).  Hence the right-hand side of (7) converges to
$C_1\rho^{a}\bigl(\int_{B_\rho}|g_\varphi|^2\bigr)^{1/2}$.

\emph{Left-hand side.}  By (6) and the Preamble item on extracting an almost
everywhere convergent subsequence from $L^2$ convergence, there is a
subsequence $\Phi_{m_l}$ with $\Phi_{m_l}\to\varphi$ almost everywhere on $K$.
The functions $|\Phi_{m_l}|^{q}$ are nonnegative and measurable and converge
almost everywhere on $K$ to $|\varphi|^{q}$, so Fatou's lemma (Preamble) gives
\[
\int_K|\varphi|^{q}\ \le\ \liminf_{l\to\infty}\int_K|\Phi_{m_l}|^{q} .
\]
Again all the integrands vanish off $K$, so these are integrals over $B_\rho$.
Combining with (7) along the subsequence and with the convergence of the
right-hand side,
\[
\Bigl(\int_{B_\rho}|\varphi|^{q}\Bigr)^{1/q}
\le C_1\rho^{a}\Bigl(\int_{B_\rho}|g_\varphi|^2\Bigr)^{1/2}.
\tag{8}
\]
(The passage from the bound on $\int|\varphi|^q$ to the bound on its $1/q$-th
power uses that $t\mapsto t^{1/q}$ is nondecreasing on $[0,\infty)$.)  Note in
particular that $\int_{B_\rho}|\varphi|^q$ is \emph{finite}: the right-hand side
of (8) is finite, since $|g_\varphi|\in L^2(B_\rho)$ by \noderef{MemW1pOn}.  So
the integral on the left of the assertion is a genuine finite integral and not
the value a divergent one would be assigned.

\emph{Step 6: from (8) to the averaged form.}  Write
$P=\bigl(\int_{B_\rho}|\varphi|^q\bigr)^{1/q}$ and
$G=\bigl(\int_{B_\rho}|g_\varphi|^2\bigr)^{1/2}$, both finite and nonnegative.
The asserted inequality is
\[
|B_\rho|^{-1/q}P\ \le\ C\rho\,|B_\rho|^{-1/2}G ,
\quad\text{i.e.}\quad
P\ \le\ C\rho\,|B_\rho|^{\frac1q-\frac12}G ,
\]
and $|B_\rho|^{\frac1q-\frac12}=(\omega_n\rho^n)^{\frac1q-\frac12}
=\omega_n^{\frac1q-\frac12}\rho^{\,n(\frac1q-\frac12)}$ by the Preamble item on
the volume of a ball.  For $n\ge3$, $q=\frac{2n}{n-2}$ gives
$n\bigl(\frac1q-\frac12\bigr)=n\bigl(\frac{n-2}{2n}-\frac12\bigr)=-1$, so the
asserted inequality is $P\le C\omega_n^{\frac1q-\frac12}G$ and (8) with $a=0$
delivers it with $C=C_1\omega_n^{\frac12-\frac1q}$.  For $n=2$, $q=4$ gives
$n\bigl(\frac1q-\frac12\bigr)=2\bigl(\frac14-\frac12\bigr)=-\tfrac12$, so the
asserted inequality is $P\le C\omega_2^{-1/4}\rho^{1/2}G$ and (8) with
$a=\tfrac12$ delivers it with $C=C_1\omega_2^{1/4}$.  In both cases $C>0$,
being a product of positive numbers, and $C$ depends only on $n$ and $C_6$.
This completes the proof.
\end{proof}
